\section{Code and experiment running snippets}
\label{sec:code}
All code, every result file, and the scripts that draw both figures are at
\url{https://github.com/WMSlalalala/verifier-supervised-lookback}. The repository README documents
each file, class and function, states which parts we wrote ourselves (all of \texttt{src/}; the
external dependencies are EasyOCR, \texttt{transformers}, \texttt{peft} and the Qwen2.5 models), and
lists the commands below. The DPO loss, the McNemar test, the BM25 scorer and the trajectory-grouped
fold splitter are hand-written rather than taken from a library.

\paragraph{Build the dataset and the labels.} \texttt{sample\_more.py} picks trajectories and downloads
their screenshots, \texttt{ocr\_shots.py} runs EasyOCR over them, and \texttt{verify\_more.py} applies
the verifier, writing both the cases and the label sets.
\begin{lstlisting}[language=bash]
MINLEN=16 MINT=12 N=400 python src/sample_more.py  # 400 trajectories
python src/ocr_shots.py          # resumable; appends to parse_odyssey.jsonl
python src/verify_more.py        # -> cases_more.json + valid_screens_more.json
python src/build_scorer_data.py  # -> 188 cases, 2,861 pairs, 5 folds
\end{lstlisting}

\paragraph{Train both instantiations and the ablations.} \texttt{FOLD} selects the held-out fold,
\texttt{TAG} names the run, and \texttt{NODIST}, \texttt{POSW} and \texttt{MODEL} drive the three
ablations. One fold takes 13--27 minutes on one 48GB GPU.
\begin{lstlisting}[language=bash]
FOLDS="0 1 2 3 4"; M3B=Qwen/Qwen2.5-3B-Instruct
for f in $FOLDS; do FOLD=$f python src/train_scorer.py;   done  # scorer
for f in $FOLDS; do FOLD=$f python src/train_fullhist.py; done  # full history
python src/score_untrained.py                      # same prompt, no adapter
for f in $FOLDS; do FOLD=$f TAG=_nodist NODIST=1 python src/train_scorer.py; done
for f in $FOLDS; do FOLD=$f TAG=_posw1  POSW=1   python src/train_scorer.py; done
for f in $FOLDS; do FOLD=$f TAG=_3B MODEL=$M3B   python src/train_scorer.py; done
\end{lstlisting}

\paragraph{Read stage, its controls, and the prompted baselines.}
\begin{lstlisting}[language=bash]
SRC=trained HIST=1 python src/read_scorer.py     # read with history
SRC=trained        python src/read_scorer.py     # single-screen read
for m in gold recent random top2; do
  SRC=trained HIST=1 MARK=$m python src/read_scorer.py            # controls
done
CASES=long MODE=B python src/eval_long.py        # full history, prompted
CASES=long MODE=D python src/eval_long.py        # screen index, prompted
\end{lstlisting}

\paragraph{External set: train once on all 188, evaluate once.}
\begin{lstlisting}[language=bash]
FOLD=ext python src/train_scorer.py
FOLD=ext python src/train_fullhist.py
SRC=trained SET=_more HIST=1 python src/read_scorer.py
\end{lstlisting}

\paragraph{Reproduce every table and figure.} \texttt{analyze\_all.py} reads the committed result
files and prints Tables~\ref{tab:main}--\ref{tab:abl} with their paired tests; it needs no GPU.
\begin{lstlisting}[language=bash]
python src/analyze_all.py
python figures/make_fig1.py && python figures/make_fig2.py
\end{lstlisting}

\paragraph{Running on shared GPUs.} The \texttt{run\_*.sh} drivers wrap the commands above. Each
waits for a GPU with enough free memory before starting a job (\texttt{pick()}), skips any stage
whose output already exists, and retries on failure, so a run that is interrupted or evicted resumes
where it stopped. This mattered: one stage of our experiments was evicted mid-run by a co-tenant job
and resumed from the completed folds.
